% ============================================================
%  Глава 9 — A* Pathfinding и Bresenham Line-of-Sight
%  Файл: compute_node/simulator.py (строки 103--280)
% ============================================================
\chapter{A* Pathfinding и алгоритм Брезенхэма}
\label{ch:pathfinding}

\begin{filebox}[title={Файл исходного кода}]
\filepath{compute\_node/simulator.py} \quad (строки 103--280, Python)
\end{filebox}

\section{Постановка задачи}

Робот должен найти оптимальный путь от текущей позиции до цели,
обходя стены арены и запретные зоны. Используется \textbf{A*} ---
классический информированный алгоритм поиска на графе.

\section{Построение сетки занятости}

\begin{filebox}[title={\filepath{simulator.py}, строки 109--137 --- \_build\_grid()}]
\begin{lstlisting}[style=pythonstyle, numbers=left, firstnumber=109]
def _build_grid(arena, zones, robot_radius):
    """Build occupancy grid: True = blocked, False = free."""
    cols = int(arena.width / PATH_GRID_RES)   # 0.05 m/cell
    rows = int(arena.height / PATH_GRID_RES)
    grid = [[False] * cols for _ in range(rows)]
    margin = robot_radius + PATH_SAFETY_MARGIN  # 0.12 + 0.10 = 0.22 m

    for r in range(rows):
        for c in range(cols):
            wx = (c + 0.5) * PATH_GRID_RES
            wy = (r + 0.5) * PATH_GRID_RES
            # Block near arena walls
            if (wx < margin or wx > arena.width - margin or
                    wy < margin or wy > arena.height - margin):
                grid[r][c] = True
                continue
            # Block inside forbidden zones (with safety margin)
            for z in zones:
                zx1, zy1 = z['x1'] - margin, z['y1'] - margin
                zx2, zy2 = z['x2'] + margin, z['y2'] + margin
                if zx1 <= wx <= zx2 and zy1 <= wy <= zy2:
                    grid[r][c] = True
                    break
    return grid, rows, cols
\end{lstlisting}
\end{filebox}

Арена дискретизируется на ячейки $\delta = 5\,\text{см}$:

\begin{equation}
  \text{cols} = \left\lfloor \frac{W}{\delta} \right\rfloor, \quad
  \text{rows} = \left\lfloor \frac{H}{\delta} \right\rfloor
  \label{eq:grid_size}
\end{equation}

Зона безопасности вокруг каждого препятствия --- \textbf{раздутие Минковского}:

\begin{equation}
  \text{margin} = r_{\text{robot}} + r_{\text{safety}} = 0.12 + 0.10 = 0.22\,\text{м}
  \label{eq:minkowski_margin}
\end{equation}

\section{Преобразование координат}

\begin{filebox}[title={\filepath{simulator.py}, строки 140--147}]
\begin{lstlisting}[style=pythonstyle, numbers=left, firstnumber=140]
def _world_to_grid(wx, wy):
    return int(wx / PATH_GRID_RES), int(wy / PATH_GRID_RES)

def _grid_to_world(c, r):
    return (c + 0.5) * PATH_GRID_RES, (r + 0.5) * PATH_GRID_RES
\end{lstlisting}
\end{filebox}

\textbf{Мир $\to$ сетка:}
\begin{equation}
  c = \left\lfloor \frac{w_x}{\delta} \right\rfloor, \quad
  r = \left\lfloor \frac{w_y}{\delta} \right\rfloor
  \label{eq:world2grid_sim}
\end{equation}

\textbf{Сетка $\to$ мир} (центр ячейки):
\begin{equation}
  w_x = (c + 0.5)\,\delta, \quad w_y = (r + 0.5)\,\delta
  \label{eq:grid2world_sim}
\end{equation}

\section{Алгоритм A*}

\begin{filebox}[title={\filepath{simulator.py}, строки 150--213 --- find\_path()}]
\begin{lstlisting}[style=pythonstyle, numbers=left, firstnumber=173]
# A* with 8-directional movement
DIRS  = [(-1,0),(1,0),(0,-1),(0,1), (-1,-1),(-1,1),(1,-1),(1,1)]
COSTS = [1.0, 1.0, 1.0, 1.0,  1.414, 1.414, 1.414, 1.414]

def heuristic(r1, c1, r2, c2):
    dr = abs(r1 - r2)
    dc = abs(c1 - c2)
    return max(dr, dc) + 0.414 * min(dr, dc)  # octile distance

open_set = [(heuristic(sr, sc, gr, gc), 0.0, sr, sc)]
g_cost = {(sr, sc): 0.0}
came_from = {}

while open_set:
    _f, g, r, c = heapq.heappop(open_set)
    if r == gr and c == gc:
        # Reconstruct path ...
        break
    for (dr, dc), cost in zip(DIRS, COSTS):
        nr, nc = r + dr, c + dc
        if 0 <= nr < rows and 0 <= nc < cols and not grid[nr][nc]:
            ng = g + cost
            if ng < g_cost.get((nr, nc), float('inf')):
                g_cost[(nr, nc)] = ng
                f = ng + heuristic(nr, nc, gr, gc)
                came_from[(nr, nc)] = (r, c)
                heapq.heappush(open_set, (f, ng, nr, nc))
\end{lstlisting}
\end{filebox}

\subsection{Оценочная функция A*}

A* расширяет вершину с минимальным $f$:

\begin{equation}
  \boxed{f(n) = g(n) + h(n)}
  \label{eq:astar_f}
\end{equation}

где $g(n)$ --- стоимость пути от старта до $n$, $h(n)$ --- эвристика от $n$ до цели.

\subsection{Октильная эвристика}

Для 8-связного графа (с диагоналями) оптимальная допустимая эвристика:

\begin{equation}
  \boxed{h(n) = \max(\Delta r, \Delta c) + (\sqrt{2} - 1)\,\min(\Delta r, \Delta c)}
  \label{eq:octile}
\end{equation}

где $\Delta r = |r_n - r_g|$, $\Delta c = |c_n - c_g|$. В коде $\sqrt{2}-1 \approx 0.414$.

\textbf{Свойства:}
\begin{itemize}
  \item \textbf{Допустимая} (admissible): $h(n) \leq h^*(n)$ --- не переоценивает.
  \item \textbf{Монотонная} (consistent): $h(n) \leq c(n,n') + h(n')$ --- гарантирует
    оптимальность без повторного посещения.
\end{itemize}

\subsection{Стоимость переходов}

\begin{equation}
  c(n, n') = \begin{cases}
    1.0   & \text{кардинальное направление (4 стороны)} \\
    \sqrt{2} \approx 1.414 & \text{диагональное направление (4 угла)}
  \end{cases}
  \label{eq:move_cost}
\end{equation}

\section{Алгоритм Брезенхэма (Line-of-Sight)}

\begin{filebox}[title={\filepath{simulator.py}, строки 234--258 --- \_line\_of\_sight()}]
\begin{lstlisting}[style=pythonstyle, numbers=left, firstnumber=234]
def _line_of_sight(x0, y0, x1, y1, grid, rows, cols):
    """Bresenham: True if straight line is obstacle-free."""
    c0, r0 = _world_to_grid(x0, y0)
    c1, r1 = _world_to_grid(x1, y1)
    dc = abs(c1 - c0)
    dr = abs(r1 - r0)
    sc = 1 if c0 < c1 else -1
    sr = 1 if r0 < r1 else -1
    err = dc - dr
    while True:
        if grid[r0][c0]:
            return False         # blocked cell
        if r0 == r1 and c0 == c1:
            break
        e2 = 2 * err
        if e2 > -dr:
            err -= dr
            c0 += sc
        if e2 < dc:
            err += dc
            r0 += sr
    return True
\end{lstlisting}
\end{filebox}

\textbf{Алгоритм Брезенхэма} растеризует отрезок прямой на целочисленной сетке:

\begin{equation}
  \text{err} = \Delta c - \Delta r, \quad
  e_2 = 2 \cdot \text{err}
  \label{eq:bresenham_err}
\end{equation}

На каждом шаге:
\begin{itemize}
  \item Если $e_2 > -\Delta r$: $\text{err} \leftarrow \text{err} - \Delta r$, сдвиг по $c$.
  \item Если $e_2 < \Delta c$: $\text{err} \leftarrow \text{err} + \Delta c$, сдвиг по $r$.
\end{itemize}

Алгоритм работает за $O(\max(\Delta c, \Delta r))$ целочисленных операций (без деления и плавающей точки).

\section{Сглаживание пути (Line-of-Sight Pruning)}

\begin{filebox}[title={\filepath{simulator.py}, строки 261--279 --- \_smooth\_path()}]
\begin{lstlisting}[style=pythonstyle, numbers=left, firstnumber=261]
def _smooth_path(path, grid, rows, cols):
    """Reduce points via line-of-sight pruning."""
    smoothed = [path[0]]
    i = 0
    while i < len(path) - 1:
        best = i + 1
        for j in range(len(path) - 1, i + 1, -1):
            if _line_of_sight(path[i][0], path[i][1],
                              path[j][0], path[j][1],
                              grid, rows, cols):
                best = j
                break
        smoothed.append(path[best])
        i = best
    return smoothed
\end{lstlisting}
\end{filebox}

Удаляет промежуточные точки, если прямая видимость (Bresenham check) свободна.
Уменьшает число путевых точек с $O(N)$ до $O(\sqrt{N})$.

\section{Геометрия A* на сетке}

\begin{center}
\begin{tikzpicture}[scale=0.6]
  % Grid
  \draw[step=1, gray!30, very thin] (0,0) grid (8,6);

  % Obstacles
  \fill[red!20] (3,2) rectangle (5,4);
  \draw[red!60, thick] (3,2) rectangle (5,4);
  \node[font=\tiny, red!80] at (4,3) {Запретная зона};

  % Margin
  \draw[red!30, dashed] (2.5,1.5) rectangle (5.5,4.5);
  \node[font=\tiny, red!40] at (4,1.2) {margin};

  % Start
  \filldraw[blue!70] (1,1) circle (0.15);
  \node[below, font=\tiny, blue!70] at (1,0.85) {Start};

  % Goal
  \filldraw[green!70!black] (7,5) circle (0.15);
  \node[below, font=\tiny, green!70!black] at (7,4.85) {Goal};

  % Path
  \draw[blue!60, very thick, ->] (1,1) -- (2,2) -- (2,4.5) -- (3,5) -- (7,5);
  \node[font=\tiny, blue!60] at (0.5,3.5) {A* путь};
\end{tikzpicture}
\end{center}

\section{Сводная таблица параметров}

\begin{center}
\begin{tabular}{lll}
\toprule
\textbf{Параметр} & \textbf{Значение} & \textbf{Описание}\\
\midrule
\texttt{PATH\_GRID\_RES}     & $0.05\,\text{м}$ & Размер ячейки сетки\\
\texttt{PATH\_SAFETY\_MARGIN} & $0.10\,\text{м}$ & Запас безопасности\\
\texttt{ROBOT\_RADIUS}        & $0.12\,\text{м}$ & Радиус робота\\
Кардинальная стоимость        & $1.0$   & 4 направления\\
Диагональная стоимость        & $\sqrt{2} \approx 1.414$ & 4 диагонали\\
\bottomrule
\end{tabular}
\end{center}
